\section{Conclusion}

We introduced Re-prompt Tax as a way to measure hidden interaction costs for
multilingual and code-switched LLM users. On a 120-item synthetic stress pilot,
three GPT-4.1-family models exhibit measurable first-turn contract failures,
especially on implicit editing-preservation tasks. A lightweight Global
Interaction Contract prompt improves first-turn global alignment and reduces
repair burden, but residual exact-preservation and script/register/locale
failures remain. The broader takeaway is methodological: global evaluation
should measure not only whether a model can eventually produce a correct answer,
but also how much repair work users must perform to make the answer usable.
